\section{Conclusiones}
\label{sec:conclusiones}
% PRISMA ítem 23d (implicaciones)

Esta revisión sistemática incluyó \nEstudios{} estudios, de los cuales \nMultiturno{} son multiturno. Las respuestas a las preguntas de investigación se limitan a lo que permiten las amenazas de la \cref{sec:amenazas}.

Respecto de PI0, el escenario multiturno es una fracción minoritaria de los estudios con vinculación de esquemas o autocorrección como contribución central: D4, una proporción descriptiva dentro del marco de elegibilidad, es \DcuatroPct~\% (IC 95~\% de \DcuatroICinf{} a \DcuatroICsup~\%), clasificado como «\DcuatroConclusion{}» por la regla registrada, sin tendencia detectable en la prueba registrada ($p$ = \DcuatroTendenciaP{}) y con dirección decreciente, no significativa, al incluir 2026 ($p$ = \tendenciaConP{}). En PI1, la concatenación del historial es la técnica de restauración del contexto más frecuente, seguida de la reescritura de la pregunta. En PI2, predomina el historial plano, y las estructuras explícitas y los grafos son minoritarios; los contrastes dentro del estudio son mayoritariamente positivos, con sesgo de notificación posible y confianza que baja al penalizarlo. En PI3, la inferencia de \en{joins} es casi siempre implícita en el modelo, aunque la comparación con un turno no es válida por la diferencia en los valores no reportados o no extraídos. En PI4, el mecanismo de autocorrección a menudo no se reporta; la aclaración con el usuario es proporcionalmente mucho más frecuente en multiturno y la autoconsistencia por votación no figura en ninguno. En PI5, el historial de evaluación varía y con frecuencia no se reporta, y las pendientes de degradación por turno reportadas son negativas en los estratos analizados, con pocas curvas.

El número de estudios del corpus que cumplen el criterio pre-especificado de reformulación, cuyas definiciones (a), (b) y (c) estaban registradas, es \nBrechaMocho{}. Los pares observados (\nMochoAB{}, \nMochoAC{} y \nMochoBC{}) se acercan a los esperados bajo independencia (\espMochoAB{}, \espMochoAC{} y \espMochoBC{}), y el triple esperado es \espMochoABC{}; el resultado es, por tanto, consistente con una baja prevalencia conjunta y no indica que la combinación se evite, sea necesaria ni eficaz. Cada componente aparece por separado. Esta afirmación es condicional a la codificación de la extracción y a los informes no recuperados (\cref{sec:dis-brecha}). El trabajo futuro derivado atiende a brechas de cobertura, ordenadas según la solidez de la evidencia: la evaluación con historial predicho y reporte por turno, un reporte más completo de la técnica y de la evaluación, la discusión de la contaminación de los \en{benchmarks}, los datos en español y multilingües, la comparación de la aclaración con el usuario frente a otras rutas, la combinación de estado estructurado con inferencia algorítmica de \en{joins} o retroalimentación de ejecución, que es poco frecuente, y la caracterización del trabajo sin entrenamiento en este escenario.
% fuente: pi15_resultados.json (brecha_M8.componentes), revision_4f.json (A2_M8_esperados, A3_tendencia_2019_2026)
